\documentclass[10pt,a4paper]{article}
\usepackage[T1]{fontenc}
\usepackage[utf8]{inputenc}
\usepackage{lmodern,microtype}
\usepackage[margin=18mm,headheight=14pt]{geometry}
\usepackage{graphicx,xcolor,booktabs,tabularx,array,amsmath,amssymb}
\usepackage{caption,titlesec,fancyhdr,enumitem,listings}
\usepackage[colorlinks=true,linkcolor=blue!45!black,urlcolor=blue!50!black]{hyperref}
\definecolor{navy}{HTML}{173650}
\definecolor{muted}{HTML}{526574}
\definecolor{light}{HTML}{EDF3F7}
\titleformat{\section}{\Large\sffamily\color{navy}}{\thesection}{.7em}{}
\titleformat{\subsection}{\large\sffamily\color{navy}}{\thesubsection}{.7em}{}
\titlespacing*{\section}{0pt}{0pt}{9pt}
\titlespacing*{\subsection}{0pt}{8pt}{5pt}
\setlength{\emergencystretch}{1em}
\setlength{\parindent}{0pt}
\setlength{\parskip}{6pt}
\setlist[itemize]{leftmargin=15pt,itemsep=3pt,topsep=3pt}
\setlist[enumerate]{leftmargin=17pt,itemsep=3pt,topsep=3pt}
\captionsetup{hypcap=false,font=small,labelfont=sf,justification=raggedright,singlelinecheck=false,skip=5pt}
\pagestyle{fancy}
\fancyhf{}
\fancyhead[L]{\small\sffamily\color{muted}Spotter | Freight Rate Assessment}
\fancyhead[R]{\small\sffamily\color{muted}Technical report}
\fancyfoot[L]{\small\color{muted}Results snapshot: 14 October 2026}
\fancyfoot[R]{\small\thepage}
\renewcommand{\headrulewidth}{.3pt}
\newcommand{\fig}[3]{\begin{center}\includegraphics[width=\linewidth,height=#2,keepaspectratio]{#1}\captionof{figure}{#3}\end{center}}
\newcommand{\code}[1]{\nolinkurl{#1}}
\newcolumntype{Y}{>{\raggedright\arraybackslash}X}
\lstset{basicstyle=\ttfamily\footnotesize,backgroundcolor=\color{light},frame=none,breaklines=true,columns=fullflexible,keepspaces=true,xleftmargin=6pt,xrightmargin=6pt}
\hypersetup{pdftitle={Freight Rate Prediction: Technical Report},pdfauthor={},pdfsubject={Spotter machine learning assessment}}
\begin{document}
{\Huge\sffamily\color{navy}Freight Rate Prediction}\par
{\Large\sffamily Technical report}\par
{\color{muted}Asgat | Spotter assessment | 15 October 2026}
\vspace{12pt}
\section*{Objective}
Predict freight rates for individual deliveries within the US using route, equipment, weight, date and market inputs. The required outputs are predictions for 12,000 unlabeled validation loads and a 31-day December price curve for a fixed Lexington--Fort Wayne load.
\subsection*{Data and evaluation}
The labeled dataset contains 48,000 loads from January through October 2025. Supplied validation covers November--December. Model selection uses three expanding temporal folds; the final 14 labeled days form a separate 2,183-row internal test.

The selected model is CatBoost trained on $\ln(\text{posted rate})$, using 52 predictors. Hyperparameters are optimized with Optuna against weighted CV MAPE, with fold weights 0.2, 0.3 and 0.5. The submission model is refitted on all labeled rows with 285 trees.
\subsection*{Main results}
\begin{center}
\begin{tabularx}{\linewidth}{Yrr}
\toprule Evaluation & Quote $\times$ distance & CatBoost \\
\midrule Weighted CV MAPE & 12.78\% & 5.1119\% \\
October test MAPE & 21.028\% & 5.5562\% \\
October test MAE & 411.32 & 115.65 \\
October test RMSE & 838.65 & 706.52 \\
\bottomrule
\end{tabularx}
\end{center}
Most test loads are predicted closely. The 38 loads with APE above 50\% represent 1.74\% of the test and contribute 61.24\% of summed percentage errors. All test rows remain included in the reported metrics.
\subsection*{Deliverables}
The saved model, preprocessing state, validation predictions and reconstructed December predictions form a reproducible inference pipeline. Both prediction files pass the assessment scorer's structural checks. December point predictions range from 813.95 to 844.65, conditional on reconstructed market and quote inputs.
\subsection*{Scope and limitations}
Validation targets are unavailable. The internal test was reserved retrospectively and is not a pristine external benchmark. The shared market curve uses known future market inputs without price targets. The December interval approximates model error and excludes uncertainty from input reconstruction.

The report presents data quality, geographic and temporal analysis, feature engineering, validation, model performance, error analysis and inference. Distributions, parameter tables and model explanations accompany their corresponding analysis sections.

\newpage
\section{Data quality and exploratory analysis}
\fig{figures/target_distributions.pdf}{.25\textheight}{The same 48,000 training targets before and after the natural logarithm. No values are clipped or removed. Validation targets are unavailable.}
The training target spans 57.22--25,533.00 USD, with median 2,030.76. The long right tail motivates modeling $\ln(\text{posted rate})$; predictions are mapped back with the exponential. Log transformation compresses the scale without deleting extreme prices.

Train and validation have no duplicate load IDs or duplicate complete rows. Distances, quote signals and training targets are positive and finite. Weight and market index require cleaning. Empty prediction-template columns are expected outputs, not missing predictors.
\begin{center}
\par\noindent\begin{tabularx}{\linewidth}{Yrr}
\toprule Issue & Train (48,000) & Validation (12,000) \\
\midrule Missing weight & 300 (0.625\%) & 165 (1.375\%) \\
Nonpositive weight & 292 (0.608\%) & 145 (1.208\%) \\
Missing market index & 374 (0.779\%) & 249 (2.075\%) \\
Duplicate load IDs & 0 & 0 \\
\bottomrule
\end{tabularx}\par
\end{center}
Nonpositive weights become missing. Weight and market medians fit only the corresponding training partition, including within CV folds. Cleaning preserves row order, IDs and supplied coordinate fields. No target outliers are deleted.

\newpage
\subsection*{Numerical feature distributions}
\fig{figures/eda_distributions.pdf}{.77\textheight}{All eight raw numerical predictors, with common bins and independently normalized distributions for train and validation. Missing values are excluded from histograms and reported in the preceding audit; nonpositive weights remain visible before cleaning.}
Distance medians are 953.3 versus 953.6 miles; quote medians are 2.056 versus 2.051. Mean market index changes from 1.083 in train to 0.927 in validation.

\newpage
\subsection*{Equipment and date distributions}
\fig{figures/equipment_date_distributions.pdf}{.60\textheight}{All equipment categories and the monthly date distribution. Percentages use each dataset's own row count. Dates show the chronological split, not a random partition.}
The next two pages show every pickup and delivery category, including validation-only cities. Together these plots cover all original predictors. \code{load_id} is a unique identifier and is excluded from modeling; its uniqueness is checked in the data-quality audit.

\newpage
\subsection*{Pickup-city distribution}
\fig{figures/pickup_distributions.pdf}{.80\textheight}{All pickup cities, ordered by pooled train and validation frequency. Each dataset is normalized separately; a missing bar denotes zero observations.}

\newpage
\subsection*{Delivery-city distribution}
\fig{figures/delivery_distributions.pdf}{.80\textheight}{All delivery cities, ordered by pooled train and validation frequency. Each dataset is normalized separately, including cities absent from training.}

\newpage
\subsection*{New cities in validation}
Each endpoint role has 64 training cities and 72 validation cities. The eight names absent from the union of training endpoints are Allentown, Charlotte, Chicago, Jackson, Knoxville, Laredo, Norfolk and San Diego.

The colored circles summarize validation restricted to known pickup and delivery cities; the daily-load comparison on the next page uses the same subset. Positions use the separate GeoNames city reference, rather than the displaced coordinates supplied with the freight records.
\fig{figures/city_composition.pdf}{.51\textheight}{Circle area represents a city's share of all training pickup and delivery endpoints. Color shows the sum of pickup and delivery percentage changes in mean daily loads: red indicates growth, blue decline. Magenta stars identify the eight cities absent from both training endpoint roles.}
The new cities are mostly within regions represented in training. Geographic features may help generalization, but this is not a guarantee of accuracy.

For each endpoint role, growth is $(\text{validation loads/day}/\text{training loads/day}-1)\times100\%$. The color adds these two percentages; it is not the percentage change of their combined count. Counts use the full calendar length of each split, including days without loads for a city.

New cities have zero training share and undefined percentage growth, so they receive fixed-size stars. Their presence is determined from the full validation set; the colored circles use validation restricted to known pickup and delivery cities. Marker areas encode training composition, not city population or geographic size.

\newpage
\subsection*{Daily load counts and sampling shift}
Train spans 304 days and validation 61. Mean sample count rises from 157.89 to 196.72 loads per day. Restricting validation to known pickup/delivery cities leaves 10,553 rows, or 173.00 per day. Novel cities explain part, but not all, of the sample-count gap.
\fig{figures/daily_load_counts.pdf}{.43\textheight}{Daily dataset row counts for train, all validation and validation restricted to known endpoint cities. Thin, faint lines show daily counts; solid bold lines show trailing 7-calendar-day means, computed separately within each split. The dotted line marks the start of validation.}
The known-city filter checks each endpoint's training role separately: both pickup and delivery must have appeared in the corresponding training column. The remaining gap is approximately 15 daily rows even after this restriction.

These are sampled records, not verified counts of the entire market. Higher sampling density alone does not prove an increase in real freight demand. The available data do not establish the cause of the remaining jump.

\newpage
\subsection*{Global and Detroit temporal signals}
The panels below compare market index, quote signal and observed price globally and for trips from/to Detroit. Validation is restricted to known endpoint cities; its actual prices are unavailable.
\fig{figures/temporal_signals.pdf}{.55\textheight}{The three signals from EDA: global daily means, then trailing 7-calendar-day load-weighted means from and to Detroit, separately within each split. Market and quote use the left axis; posted rate uses the right. Quotes are dashed; faint Detroit points are individual market observations. Validation has no posted-rate curve.}
The temporal panels expose market regimes and differences in the signals without establishing causation. The analysis does not show that new cities cause the market-index shift.

\newpage
\subsection*{Coordinate quality and geographic assumptions}
The dataset supplies one coordinate pair per named city, constant across endpoint roles and datasets. It does not reveal individual warehouse locations within a city. Against matched official city boundaries, 70 of 72 points lie outside their named municipality; 52 are more than 100 km outside.
\fig{figures/city_boundaries.pdf}{.42\textheight}{Original supplied points and US Census named-city boundaries. Outside the named city does not mean outside every municipality, nor prove that a coordinate is incorrect.}
Original coordinates remain model inputs. Road features use a separate GeoNames city reference with a documented Washington override: routing displaced supplied points would answer a different geographic question. References are not claimed to be real warehouse sites.
References and directed route tables are saved with preprocessing. Inference needs no live routing API. Ambiguous city names use selected reference records; geographic matching remains an assumption. External data provenance is listed in the sources section.

\newpage
\section{Feature engineering}
\subsection*{Road distance and driving time}
Road distance and estimated driving time provide additional route information for freight pricing. These inputs are obtained using HERE Matrix Routing, with cached route geometries used for visual checks.
\fig{figures/routes.pdf}{.38\textheight}{Example truck routes between the separate city reference points: Dallas--Houston, Detroit--Atlanta and Los Angeles--New York. Circles mark pickup; triangles mark delivery.}
The features include road distance, driving time and ratios comparing road, straight-line and supplied distance. Original supplied coordinates remain separate model inputs.

These are routing estimates, not verified historical 2025 journeys or confirmed warehouse-to-warehouse paths. They do not include loading or driver rest. The geographic reference and routing assumptions are saved with preprocessing.

\newpage
\subsection*{Market-index decomposition}
The observed daily market mean is fitted by up to four chronological regimes:
\[
f_k(t)=a_k+b_kt+c_kt^2+u_k\sin(\theta)+v_k\cos(\theta),\qquad
 t=\frac{\text{days since 2025-01-01}}{365.25},\quad \theta=\frac{2\pi\,\text{weekday}}7.
\]
Breakpoints minimize daily squared error; BIC selects the number of segments. A minimum of seven observed days permits a new short regime without forcing it to absorb the previous one. Global daily-mean fit RMSE is 0.002588 index units.
\fig{figures/market_decomposition.pdf}{.48\textheight}{Fitted expectation, quadratic trend and weekly harmonic. The dotted date marks supplied validation. Its known market inputs participate in this descriptive fit.}
\begin{center}\small
\begin{tabular}{lrrrrr}
\toprule Regime dates & $a$ & $b$ & $c$ & $u$ & $v$ \\
\midrule Jan 1--Apr 15 & 0.88003 & 0.96751 & 0.05180 & 0.05497 & -0.08454 \\
Apr 16--Jul 24 & 0.81194 & 2.48109 & -3.15744 & 0.05449 & -0.08337 \\
Jul 25--Oct 2 & 1.79519 & -1.56563 & 0.41259 & 0.05485 & -0.08427 \\
Oct 3--Dec 31 & 3.47174 & -5.71274 & 3.19806 & 0.05439 & -0.08352 \\
\bottomrule
\end{tabular}\end{center}
The equation uses known train plus validation \emph{inputs}, without price targets. Residual reference distributions and equipment/lane residual statistics fit each training partition only. This covariate-informed batch setup is not a strict simulation of online forecasting without future market inputs; deployment would require a separate input-availability policy.

\newpage
\subsection*{Predictor inventory and preprocessing}
The model uses 52 features, including three categorical columns: pickup, delivery and equipment. Load IDs, raw date metadata, targets and output fields are excluded.
\par\noindent\begin{tabularx}{\linewidth}{lrY}
\toprule Group & Count & Information \\
\midrule Original predictors & 11 & Cities, equipment, supplied coordinates, distance, weight, market and quote. \\
Cyclical time & 7 & Sine/cosine of weekday, within-month and within-year phase; weekend flag. \\
Road routes & 5 & HERE miles/hours, road-to-air and input-to-road ratios, availability. \\
Market decomposition & 16 & Trend, harmonic, expectation, derivatives, residual statistics and five regime coefficients. \\
US holidays & 7 & Actual/observed flags, distance to adjacent holidays and nearby-day indicators. \\
Historical city inputs & 6 & Endpoint shares, directional balance and destination/origin diversity. \\
\bottomrule
\end{tabularx}\par
City aggregates use previous days only for training; the entire current day is excluded before updates. Test and validation use frozen training history. No target encoding or price-derived aggregate enters the predictors. Calendar holidays are known in advance; there is no raw elapsed-date feature.
Route distances and durations augment the supplied load inputs without replacing their original coordinate columns. Correlated distance features are interpreted jointly in the SHAP analysis later in the report.

\newpage
\section{Chronological validation}
The final 14 labeled days, October 18--31, form an internal test of 2,183 loads. Development contains 45,817 earlier rows. Expanding folds split unique dates, keeping a day's loads together.
\fig{figures/time-series-cv.png}{.28\textheight}{Three development folds and the separate holdout. There is no random mixing of days across a train/validation boundary.}
\begin{center}\small
\begin{tabular}{cllrr}
\toprule Fold & Training dates & Validation dates & Train & Valid \\
\midrule 1 & Jan 1--Mar 15 & Mar 16--May 26 & 11,673 & 11,518 \\
2 & Jan 1--May 26 & May 27--Aug 6 & 23,191 & 11,499 \\
3 & Jan 1--Aug 6 & Aug 7--Oct 17 & 34,690 & 11,127 \\
\bottomrule
\end{tabular}\end{center}
Imputation, market-residual distributions and city histories fit each fold's training rows. The shared market equation is the disclosed exception: it uses known future covariates, never price targets.

Optuna minimizes original-scale MAPE with recency weights:
\[
J=0.2\,\mathrm{MAPE}_1+0.3\,\mathrm{MAPE}_2+0.5\,\mathrm{MAPE}_3.
\]
Here $\mathrm{MAPE}=\frac{100}{n}\sum_i |y_i-\hat y_i|/y_i$; a value of 1\% means mean relative error, not classification accuracy. No pooled OOF predictions or score are computed. These windows also choose parameters and early stopping, so CV scores are development estimates after selection.

The test model fits development only, using the CV-selected fixed tree count; its labels are read for evaluation afterward. A separate submission model then refits all 48,000 labeled rows. That full-train model predicts supplied validation, not the internal test.

The holdout dates appeared in earlier experiments and were reserved retrospectively. Its score is informative, but not equivalent to a pristine one-time external evaluation. Fixed boundaries and this disclosure are retained throughout the report.

\newpage
\section{Modeling and hyperparameter optimization}
\subsection*{Product baseline}
The initial baseline predicts $\hat y=\text{quote\_signal}\times\text{distance}$, with no fitted multiplier. It gives 21.028\% MAPE on the complete October test and provides a reference for the learned model.
\subsection*{CatBoost and Optuna}
The active Optuna study has 50 completed TPE trials, seed 42. Each trial trains on $\ln(y)$ with RMSE loss and log-RMSE early stopping after 100 rounds, up to 1,500 trees. Selection uses weighted original-scale MAPE.

Search covers depth 4--8; learning rate 0.02--0.15; L2 0.1--30; random strength 0.001--3; column sampling 0.4--1; borders 64/128/254; SymmetricTree, Depthwise and Lossguide; Newton/Gradient leaf estimation; 1--10 leaf iterations; and No/AnyImprovement backtracking. Leaf-size and leaf-count parameters are conditional on growth policy.
\begin{center}
\par\noindent\begin{tabularx}{\linewidth}{YrYr}
\toprule Parameter & Value & Parameter & Value \\
\midrule Depth & 6 & Learning rate & 0.0569703 \\
L2 leaf regularization & 2.28909 & Random strength & 1.16862 \\
Column sampling & 0.977059 & Global borders & 254 \\
Growth policy & \multicolumn{1}{l}{SymmetricTree} & Leaf method & Newton \\
Leaf iterations & 1 & Final trees & 285 \\
\bottomrule
\end{tabularx}\par
\end{center}
Backtracking is AnyImprovement. Fold stopping counts are 218, 285 and 427; their median, 285, fixes test and submission tree counts. Active-study trial 41 is selected. SQLite study names fingerprint prepared fold values and configuration.
\fig{figures/distance_borders.pdf}{.26\textheight}{Extra quantization search for golden feature \code{distance}, with other settings fixed. The current pipeline selects 254; additional bins do not improve it.}
\code{golden_features} accepts numeric column names and selects a separate border count for each, sequentially under the same objective. Only results from the current feature/data fingerprint are reported; earlier pipeline experiments are not mixed into this comparison.

\newpage
\section{Model performance}
The direct baseline is $\hat y=\text{quote\_signal}\times\text{distance}$, without a fitted multiplier or correction. It gives 21.028\% MAPE on the October test.

CatBoost is trained on the natural logarithm of price, with hyperparameters tuned using Optuna. Selection uses the weighted three-fold MAPE; parameters, metrics and the final model are saved. The selected model gives 5.1119\% weighted CV MAPE and 5.5562\% test MAPE. Both methods are evaluated on identical rows.
\fig{figures/catboost-cv-metrics.png}{.25\textheight}{CatBoost and the product baseline: MAE, RMSE and MAPE improve on every split.}
\begin{center}\small
\begin{tabular}{lrrrr}
\toprule Split & CatBoost MAPE & Product MAPE & CatBoost MAE & Product MAE \\
\midrule Fold 1 & 6.200\% & 16.298\% & 133.387 & 323.250 \\
Fold 2 & 5.006\% & 12.881\% & 127.400 & 268.881 \\
Fold 3 & 4.740\% & 11.310\% & 108.474 & 239.234 \\
October test & 5.556\% & 21.028\% & 115.648 & 411.320 \\
\bottomrule
\end{tabular}\end{center}
\fig{figures/test-quote-product-baseline.png}{.30\textheight}{Individual predictions and daily means on the test. Similar daily mean rates do not imply equal load-level accuracy: errors can cancel during averaging.}
Test median APE is 1.897\%, P90 APE is 4.624\%, and 97.939\% of loads are within 10\% of observed price. RMSE remains sensitive to large deviations. The next pages examine these cases without excluding them from headline metrics.

\newpage
\section{Error analysis}
Residual is $y-\hat y$: positive means underprediction. All test diagnostics use the development-only model, not the full-train submission model.
\fig{figures/test-residuals.png}{.50\textheight}{Signed percentage errors, residual magnitude, actual versus predicted price and daily errors. Log and symmetric-log axes expose extremes without deleting rows.}
Most predictions track actual rates closely, but extreme errors occur in both directions. Price-scale errors tend to grow with price; relative metrics therefore complement MAE and RMSE. Daily spikes may reflect a few loads rather than a uniform bias.
\begin{center}
\par\noindent\begin{tabularx}{\linewidth}{Yrrr}
\toprule Test group & Rows & Mean APE & Total APE share \\
\midrule APE at most 10\% & 2,138 & 2.163\% & 38.131\% \\
APE 10--50\% & 7 & 10.945\% & 0.632\% \\
Severe overprediction & 22 & 286.162\% & 51.904\% \\
Severe underprediction & 16 & 70.756\% & 9.334\% \\
\bottomrule
\end{tabularx}\par
\end{center}
The 38 loads above 50\% APE are 1.74\% of rows but contribute 61.24\% of summed percentage errors. The 2.163\% ordinary-case metric is diagnostic, not an alternative overall score. Removing rows based on their test errors would bias evaluation; overall MAPE stays 5.556\%.

The extreme prices also differ sharply from similar historical loads. Missing pricing information or target anomalies are plausible explanations, not established causes. The following analysis compares these cases with historical peers.

\newpage
\subsection*{Characteristics of the largest errors}
All severe-error lanes were observed in development. Median historical lane frequency is about 14, comparable with ordinary cases. Weight, quote and route ratios overlap the rest of the data; unseen geography does not explain these particular errors.
\fig{figures/test-outlier-comparison.png}{.29\textheight}{Severe-error targets separate above and below the ordinary price relationship, while predictions stay closer to it. Target-derived ratios are diagnostics only.}
For each of the five worst cases, 20 approximate historical peers share equipment and are nearest in standardized distance, weight, quote and market inputs. Matching never uses targets or prediction errors; peers are not necessarily the same lane.
\begin{center}\small
\begin{tabular}{lrrrr}
\toprule Load ID & Actual & Predicted & Peer median & APE \\
\midrule TR-046400 & 123.32 & 657.79 & 852.36 & 433.40\% \\
TR-046179 & 643.58 & 3,292.46 & 2,970.97 & 411.59\% \\
TR-046174 & 880.13 & 4,404.34 & 4,282.88 & 400.42\% \\
TR-046260 & 352.82 & 1,677.18 & 1,617.79 & 375.36\% \\
TR-047504 & 880.04 & 4,167.31 & 4,185.07 & 373.54\% \\
\bottomrule
\end{tabular}\end{center}
These actual prices are 14--22\% of approximate peer medians. Median actual/predicted is 0.248 for severe overprediction and 3.225 for severe underprediction. Their median quote signals, roughly 1.94--2.02, resemble ordinary cases at 1.99.

An unobserved pricing factor or target anomaly is plausible, not established. Small city/equipment enrichments do not prove a reliable input-based detection rule. The analysis neither establishes an irreducible error floor nor proves that 1\% MAPE is mathematically impossible.

\newpage
\section{SHAP interpretation of the final model}
Native CatBoost contributions explain raw model output:
\[\ln(\hat y_i)=\phi_{0,i}+\sum_j\phi_{ij}.\]
Positive values raise predicted log rate; negative values lower it. Explanations pass numerical additivity checks. Feature ranking uses all explained rows; points show a deterministic sample of up to 1,500 loads and the top 20 features.
\fig{figures/shap-validation.png}{.62\textheight}{Full-train model explanations on supplied validation. Numeric colors indicate low-to-high values; categorical points are gray. No validation targets are needed.}
Distance dominates the explanations, while road distance and driving time also contribute. SHAP describes model behavior; it does not prove that adding a feature improves held-out accuracy. Correlated features should be read jointly. A small quote contribution does not prove that quote is useless or absent from interactions. Fold-specific explanations are shown next.

\newpage
\subsection*{Fold 1: held-out SHAP values}
Training: January 1--March 15; validation: March 16--May 26. Imputation and residual statistics fit this fold's training rows; the market equation uses known future covariates. Fold MAPE is 6.200\%, with 218 trees.
\fig{figures/shap-fold_1.png}{.69\textheight}{Fold 1 explanations on its validation rows, in log-rate units.}

\newpage
\subsection*{Fold 2: held-out SHAP values}
Training: January 1--May 26; validation: May 27--August 6. Fold MAPE is 5.006\%, with 285 trees. Explanations use the fold model, not the submission model.
\fig{figures/shap-fold_2.png}{.69\textheight}{Fold 2 validation explanations; ranking uses all rows and point sampling is deterministic.}

\newpage
\subsection*{Fold 3: held-out SHAP values}
Training: January 1--August 6; validation: August 7--October 17. Fold MAPE is 4.740\%, with 427 trees. The most recent fold receives weight 0.5 in the optimization objective.
\fig{figures/shap-fold_3.png}{.69\textheight}{Fold 3 held-out explanations. Differences in attribution do not establish causal changes.}

\newpage
\subsection*{Decision paths for the five worst test loads}
Paths explain the development-only model, accumulating contributions from expected log rate to each prediction. Features outside the top 20 are summed into one row, preserving additivity.
\fig{figures/test-worst-five-shap.png}{.72\textheight}{Top circles mark predicted log prices and crosses mark actual log prices. Horizontal gaps are errors, not SHAP contributions. Targets are shown only for comparison.}
The model reacts to route length, but these actual prices are far lower than predicted. SHAP explains model behavior; it does not establish why the observed price differs or prove that a target is wrong.

\newpage
\section{Inference and prediction outputs}
\subsection*{Validation predictions}
All 12,000 validation loads use the full-train model. Actual prices are unavailable. These graphs describe predictions only and do not measure accuracy.
\fig{figures/inference-predictions.png}{.49\textheight}{Validation predicted-price distribution, daily mean and P10--P90 across loads, distance dependence and equipment groups. Across-load spread is not an interval for one load.}

\newpage
\subsection*{December input reconstruction}
The prescribed load is Lexington to Fort Wayne, 360 miles, Dry Van, 32,000 lb. Dates are supplied; coordinates, quote and market are absent. Unique source city coordinates are Lexington $(36.99152,-84.99876)$ and Fort Wayne $(41.31561,-85.36206)$. Routing references remain separate.

Validation has 12 loads on this directed lane, eight in December. Three December loads also use Dry Van. Their signal values are copied on those dates. They have different weights/distances from the prescribed load, so the transferred signals are proxies for the fixed-load scenario.
\begin{center}\small
\begin{tabular}{lrrl}
\toprule Date & Market index & Quote signal & Source load ID \\
\midrule Dec 5 & 0.97834 & 1.92127 & TE-006705 \\
Dec 9 & 0.92830 & 2.10398 & TE-007537 \\
Dec 21 & 0.81167 & 2.33895 & TE-009851 \\
\bottomrule
\end{tabular}\end{center}
\fig{figures/december-quote-signal-fit.png}{.23\textheight}{Missing quotes use $q(d)=1.859006+0.0245994d$, with $d$ days since December 1. Observations stay unchanged; values outside their date span are extrapolated.}
\fig{figures/december-market-index-fit.png}{.23\textheight}{Missing market values use the existing equation plus $\delta=-0.006200$, the mean observed residual. Three-point fit RMSE is 0.01910; original values are preserved.}
No price target enters reconstruction. Only three matched observations are available, so this is an uncertain input scenario, not a separately validated lane-level forecast.

\newpage
\subsection*{December predictions and uncertainty}
All 52 model features are constructed from completed inputs and saved state. Inference does not refit the model or request routes. Point prices span 813.95--844.65, with mean 825.12.
\fig{figures/candidate_december.png}{.26\textheight}{The required chart, generated by the assessment scorer from the completed seven-column December file. Prescribed raw load attributes are fixed; only date varies.}
\fig{figures/december-predictions-interval.png}{.29\textheight}{Approximate central 90\% model-error band transferred from the October test. This is not a guaranteed conditional prediction interval.}
The 5th/95th empirical quantiles of actual/predicted on 2,183 holdout loads give
\[[\hat yQ_{0.05},\hat yQ_{0.95}]=[0.959735\hat y,1.051664\hat y].\]
Coverage on that calibration sample is 89.922\%; this is not an independent coverage test. Transfer across models, periods and lanes is an assumption.

The band excludes uncertainty in reconstructed market/quote values, the base market curve and geographic references. Input-uncertainty Monte Carlo propagation was discussed but not implemented. This band must not be labeled total uncertainty or a calibrated 90\% guarantee for the fixed lane.

The weekly pattern is a model prediction conditional on these inputs, not a guarantee that certain weekdays are always cheaper. Next steps are to investigate unusual prices, obtain better shipment context and validate improvements on a fresh future period.

\newpage
\section{Reproducibility, outputs and limitations}
\begin{enumerate}
\item \code{notebooks/eda.ipynb}: quality, distributions, temporal comparisons and geographic audit.
\item \code{notebooks/modeling.ipynb}: features, folds, tuning, metrics, SHAP, residuals and artifact saving.
\item \code{notebooks/inference.ipynb}: artifact loading, validation predictions, December reconstruction and predictions, scorer.
\end{enumerate}
Shared helpers are in \code{notebooks/utils.py}; reusable \code{TimeSeriesCV} and \code{HParamsOpt} are in \code{notebooks/templates.py}. \code{scripts/build_report.py} rebuilds report figures and evidence without training.
\subsection*{Artifacts under \code{outputs/}}
\par\noindent\begin{tabularx}{\linewidth}{lY}
\toprule Artifact & Purpose \\
\midrule \code{final-model/model.cbm} & Native full-train model. \\
\code{final-model/preprocessing.pkl} & Feature state/order, categories and routing references. \\
\code{final-model/metadata.json} & Parameters, versions, dates, metrics and interval calibration. \\
\code{validation-predictions.csv} & 12,000 unique IDs and positive finite rates. \\
\code{december-predictions.csv} & 31 fixed-route rows in the original seven-column schema. \\
\code{december-prediction-intervals.csv} & Separate lower and upper estimates. \\
\code{scorer_results/candidate_december.png} & Required December chart. \\
\bottomrule
\end{tabularx}\par
\subsection*{Run from the repository root}
\begin{lstlisting}
python -m pip install -r requirements.txt
# Execute modeling.ipynb, then inference.ipynb in order.
python notebooks/score.py \
  --predictions outputs/validation-predictions.csv \
  --december-predictions outputs/december-predictions.csv \
  --output-dir outputs/scorer_results
python -m unittest discover -s tests
python scripts/build_report.py
cd output/pdf
pdflatex -interaction=nonstopmode freight_rate_report.tex
\end{lstlisting}
Use a notebook-capable environment; the verified runtime includes Jupyter/nbclient. Native-model reload checks preserve predictions. Saved versions: CatBoost 1.2.10, NumPy 2.1.3, pandas 2.3.3, scikit-learn 1.8.0. Evidence JSON includes source-artifact SHA-256 hashes.

The scorer checks file contracts, not hidden validation accuracy. CV reuses selection windows; the holdout is retrospective; future market covariates are used; December inputs are reconstructed. These limits accompany the scores. The repository contains the report and reproducible artifacts; the main report presents the modeling and inference methodology.

\newpage
\section{Sources and evidence}
\begin{itemize}
\item Assessment: local \code{data/README.md} and \code{notebooks/score.py}. The referenced assessment PDF is absent from the workspace.
\item Results: executed notebooks, saved model metadata and current golden-feature results; \code{output/pdf/evidence.json} records evidence and hashes.
\item City references: \href{https://download.geonames.org/export/dump/}{GeoNames cities15000}, CC BY 4.0; records in \code{data/city-coordinates.csv}.
\item Polygons: \href{https://www.census.gov/geographies/mapping-files/time-series/geo/tiger-line-file.html}{US Census TIGER/Line}, 2025 assets; local build/audit logic is in \code{scripts/build_city_boundaries.py}.
\item Routes: cached HERE truck estimates, snapshot \code{2026-10-12-city-routing-v2}. No credentials appear in report artifacts.
\item Contribution semantics and additivity: \href{https://catboost.ai/docs/en/concepts/shap-values}{CatBoost ShapValues documentation}. Native values are plotted with Matplotlib.
\end{itemize}
\end{document}
